\documentclass[11pt]{article}
\usepackage{amsmath,amssymb,geometry,booktabs,xcolor}
\geometry{margin=2.5cm}
\title{Exact Koide Chain from $\delta=15^\circ$}
\author{Notebook -- phys6gpd}
\date{}

\begin{document}
\maketitle

\section{Tuple 1: $(0,\, q_2^2,\, q_3^2)$ with $\delta = 15^\circ$}

The Koide parametrization writes the three charges (square roots of masses) as
\[
  q_k = \sqrt{M_0}\!\left(1 + \sqrt{2}\,\cos\!\left(\tfrac{2\pi k}{3} + \delta\right)\right),
  \qquad k=1,2,3.
\]
Setting $\delta = \pi/12 = 15^\circ$, the $k\!=\!1$ entry vanishes:
\[
  \cos\!\left(\tfrac{2\pi}{3}+\tfrac{\pi}{12}\right)
  = \cos\tfrac{3\pi}{4} = -\tfrac{1}{\sqrt{2}},
\]
so $q_1 = \sqrt{M_0}(1-1)=0$.  The surviving charges are
\begin{align}
  q_2 &= \sqrt{M_0}\!\left(1+\sqrt{2}\cos\tfrac{17\pi}{12}\right)
       = \sqrt{M_0}\,\frac{3-\sqrt{3}}{2}, \\[4pt]
  q_3 &= \sqrt{M_0}\!\left(1+\sqrt{2}\cos\tfrac{\pi}{12}\right)
       = \sqrt{M_0}\,\frac{3+\sqrt{3}}{2}.
\end{align}
The unnormalised masses are $q_k^2$; their sum is
\[
  q_2^2 + q_3^2 = M_0\!\left(\frac{(3-\sqrt3)^2+(3+\sqrt3)^2}{4}\right)
  = M_0\,\frac{24}{4} = 6\,M_0.
\]
Setting $\sqrt{M_0}=1$ henceforth, the Koide scale parameter is
$Q_0^{(1)} = (0 + q_2^2 + q_3^2)/3 = 2$.

Normalising so that $0 + x^2 + y^2 = 1$ (with $x < y$):
\[
  \boxed{x^2 = \frac{2-\sqrt{3}}{4}\approx 0.0670,
  \qquad
  y^2 = \frac{2+\sqrt{3}}{4}\approx 0.9330.}
\]
These correspond to $m_s$ and $m_c$ in the $(u,s,c)$ triplet of the Koide chain,
with $m_u=0$.

\section{Tuple 2: $(-q_2,\, q_3,\, z)$ --- predicting the $b$ charge}

The next link in the chain takes charges $(-q_2,\,q_3,\,z)$ with $z>0$,
corresponding to the $(s,c,b)$ triplet with the sign flip $-\!\sqrt{m_s}$.
The Koide condition reads
\[
  \frac{(-q_2+q_3+z)^2}{q_2^2+q_3^2+z^2} = \frac{3}{2}.
\]

\subsection{Key intermediate: $q_3 - q_2 = \sqrt{3}$}

\[
  q_3 - q_2 = \frac{(3+\sqrt3)-(3-\sqrt3)}{2} = \sqrt{3},
  \qquad q_2^2 + q_3^2 = 6.
\]

\subsection{Solving the quadratic}

The Koide equation becomes
\[
  (\sqrt{3}+z)^2 = \tfrac{3}{2}(6+z^2),
\]
\[
  3 + 2\sqrt{3}\,z + z^2 = 9 + \tfrac{3}{2}z^2,
\]
\[
  \tfrac{1}{2}z^2 - 2\sqrt{3}\,z + 6 = 0
  \qquad\Longrightarrow\qquad
  z^2 - 4\sqrt{3}\,z + 12 = 0.
\]
The discriminant is $\Delta = 48 - 48 = 0$: a \textbf{double root},
\[
  \boxed{z = 2\sqrt{3}, \qquad z^2 = m_b = 12.}
\]

\subsection{Verification and Koide angle}

\[
  Q = \frac{(\sqrt{3}+2\sqrt{3})^2}{6 + 12}
  = \frac{27}{18} = \frac{3}{2}.\;\checkmark
\]

The scale parameter is $Q_0^{(2)} = (6+12)/3 = 6 = 3\,Q_0^{(1)}$.
The Koide angle evaluates to $\delta_2 = 45^\circ = 3\times 15^\circ$
in the same phase convention, confirming
$\delta_{scb} = 3\,\delta_{0sc}$ at the exact chain point.

\subsection{Mass ratios}

\[
  m_s : m_c : m_b
  = \tfrac{(3-\sqrt3)^2}{4} : \tfrac{(3+\sqrt3)^2}{4} : 12
  = (2-\sqrt{3}) : (2+\sqrt{3}) : 8.
\]

\section{Charges vs.\ masses: when does a Koide continuation exist?}

\subsection{The general discriminant}

Given two charges $a$ and $b$,
we seek a third charge $c$ satisfying
$(a+b+c)^2 = \tfrac{3}{2}(a^2+b^2+c^2)$.
This is a quadratic in $c$ with discriminant
\[
  \boxed{\Delta = 12\bigl(a^2 + 4ab + b^2\bigr).}
\]

\subsection{Same-sign charges always continue}

When $a$ and $b$ have the same sign, $4ab > 0$,
so $\Delta > 0$: there are always two real solutions for $c$.

\subsection{Opposite-sign charges have a forbidden band}

When $a$ and $b$ have opposite signs, write $a/b = -r$ with $r>0$.
Then $\Delta = 12\,b^2(r^2 - 4r + 1)$.
The roots of $r^2 - 4r + 1 = 0$ are $r = 2 \pm \sqrt{3}$, so:
\begin{itemize}
  \item $\Delta > 0$ (two solutions) when
    $r < 2-\sqrt{3}$ or $r > 2+\sqrt{3}$;
  \item $\Delta = 0$ (double root) when $r = 2 \pm \sqrt{3}$;
  \item $\Delta < 0$ (no real solutions) when
    $2-\sqrt{3} < r < 2+\sqrt{3}$.
\end{itemize}
In terms of the mass ratio $m_1/m_2 = r^2$, the forbidden band is
$(7-4\sqrt3) < m_1/m_2 < (7+4\sqrt3)$.\footnote{Throughout
this work, masses are obtained from charges by squaring,
$m = q^2$, preserving holomorphy in $q$ as a complex variable.
An alternative construction using $m = |q|^2$ (i.e.\ $q\bar q$, where
the mass arises from phase cancellation between a charge and its conjugate)
would break holomorphy and is left for future investigation.}

\subsection{Masses vs.\ charges}

Given two \emph{masses}, one can always choose the same-sign
pattern, which always yields two Koide continuations.
\textbf{A pair of masses always has Koide descendants.}
But a pair of \emph{charges} with signs already fixed may not:
the opposite-sign pattern fails in the forbidden band.

\subsection{$\delta = 15^\circ$ sits exactly at the boundary}

At $\delta = 15^\circ$, the charge ratio $q_2/q_3 = 2 - \sqrt{3}$
is \emph{exactly} the lower boundary of the forbidden band.
This is why Tuple~2 has a double root.

\section{Perturbation: $\delta = \pi/12 + \varepsilon$}

\subsection{Perturbed charges of Tuple 1}

To first order in $\varepsilon$:
\begin{align}
  q_1 &= -\varepsilon + O(\varepsilon^2), \\
  q_2 &= \tfrac{3-\sqrt3}{2} + \tfrac{1+\sqrt3}{2}\,\varepsilon + O(\varepsilon^2), \\
  q_3 &= \tfrac{3+\sqrt3}{2} + \tfrac{1-\sqrt3}{2}\,\varepsilon + O(\varepsilon^2).
\end{align}
For $\varepsilon\neq 0$, the ``up quark'' acquires a
small mass $m_u = \varepsilon^2 + O(\varepsilon^3)$.

\subsection{Discriminant of Tuple 2}

The charge ratio $q_2/q_3$ moves away from the critical value $2-\sqrt{3}$:
$q_2/q_3 = (2-\sqrt3)\!\left(1 + \frac{4\sqrt3}{3}\,\varepsilon + O(\varepsilon^2)\right)$.
The discriminant of the Koide quadratic for Tuple~2 evaluates to
\[
  \Delta = 3\,\varepsilon\,(\varepsilon - 12) + O(\varepsilon^3).
\]
\begin{itemize}
  \item $\varepsilon = 0$: $\Delta = 0$, double root $z = 2\sqrt{3}$.
  \item $\varepsilon < 0$ (continuation-preserving side):
    $\Delta \approx 36\,|\varepsilon| > 0$, two solutions
    $z = 2\sqrt{3} \pm 6\sqrt{|\varepsilon|} + O(\varepsilon)$.
  \item $\varepsilon > 0$: $\Delta < 0$,
    no real continuation with the $(-,+)$ sign pattern.
\end{itemize}

This local bifurcation analysis concerns the exact opposite-sign continuation
of Tuple~2.  The measured positive-sign $(0,s,c)$ tuple discussed later sits
slightly on the $\varepsilon>0$ side of the critical point, whereas the
continuation branches used in the inverse-tuple analysis are tracked from the
$\varepsilon<0$ side where the real opposite-sign continuation still exists.

\section{Inverse-mass Koide tuples: $K(1/m_d,\,1/m_s,\,1/m_b)$}

We now ask: given the known masses $m_s$ and $m_b$ from the chain,
what values of $m_d$ make the \emph{inverse-mass} triplet
$(1/m_d,\,1/m_s,\,1/m_b)$ satisfy the Koide equation?

Define the \textbf{inverse charges}
\[
  w_x \equiv \frac{1}{q_x} = \frac{1}{\sqrt{m_x}},
\]
so that the ``masses'' in the inverse tuple are $w_x^2 = 1/m_x$.
With our exact values:
\[
  w_s = \frac{1}{q_2} = \frac{2}{3-\sqrt3} = \frac{3+\sqrt3}{3},
  \qquad
  w_b = \frac{1}{2\sqrt3} = \frac{\sqrt3}{6}.
\]

\subsection{Sign patterns}

The charge ratio $w_s/w_b = q_b/q_s = 4\sqrt{3}/(3-\sqrt3) = 2(1+\sqrt3) \approx 5.46$,
which exceeds $2+\sqrt3\approx 3.73$, so the opposite-sign pair $(-w_s,\,w_b)$
lies \emph{outside} the forbidden band. Both sign patterns yield real solutions:

\medskip\noindent
\textbf{Case~I:} same-sign $(+w_s,\,+w_b,\,w_d)$.\\
\textbf{Case~II:} opposite-sign $(-w_s,\,+w_b,\,w_d)$.

\subsection{Case I: $(+w_s,\,+w_b,\,w_d)$}

The quadratic in $w_d$ is
\[
  w_d^2 - 2(2+\sqrt3)\,w_d + \tfrac{3}{4} = 0,
\]
with discriminant $\Delta_{\rm I} = 25 + 16\sqrt{3}$
(not a perfect square).  The solutions are
\[
  w_d = (2+\sqrt3) \pm \tfrac{1}{2}\sqrt{25+16\sqrt3}\,,
\]
giving $m_d = 1/w_d^2$.

\subsection{Case II: $(-w_s,\,+w_b,\,w_d)$}

The quadratic in $w_d$ becomes
\[
  w_d^2 + \tfrac{2(6+\sqrt3)}{3}\,w_d + \tfrac{25+16\sqrt3}{12} = 0.
\]
The discriminant simplifies dramatically:
\[
  \boxed{\Delta_{\rm II} = 9, \qquad \sqrt{\Delta_{\rm II}} = 3.}
\]
The two solutions are
\[
  w_d = -\frac{3+2\sqrt3}{6}
  \qquad\text{and}\qquad
  w_d = -\frac{21+2\sqrt3}{6}\,,
\]
both negative (consistent with the opposite-sign pattern
$(-w_s,\,+w_b,\,w_d)$ having $w_d < 0$).

\subsubsection{Solution II-1: $w_d = -(3+2\sqrt3)/6$}

\[
  w_d^2 = \frac{(3+2\sqrt3)^2}{36} = \frac{7+4\sqrt3}{12},
  \qquad
  m_d = \frac{12}{7+4\sqrt3} = 12(7-4\sqrt3) = 12(2-\sqrt3)^2.
\]
The inverse-mass Koide angle of this solution is
$\delta = 15^\circ$ \textbf{exactly}, matching Tuple~1.

\subsubsection{Solution II-2: $w_d = -(21+2\sqrt3)/6$}

\[
  w_d^2 = \frac{(21+2\sqrt3)^2}{36} = \frac{151+28\sqrt3}{12},
  \qquad
  m_d = \frac{12}{151+28\sqrt3} = \frac{12(151-28\sqrt3)}{143^2}.
\]

\subsection{Summary table}

Define $W_0 = (w_d^2+w_s^2+w_b^2)/3$ for the inverse tuple and
recall $Q_0 = Q_0^{(2)} = 6$ for the direct $(s,c,b)$ tuple.

\newcommand{\dec}[1]{{\scriptscriptstyle\color{gray}(#1)}}

\begin{table}[h]
\centering
\small
\begin{tabular}{@{}l c c c c c c@{}}
\toprule
Case & $w_d$ & $w_d^2=\frac{1}{m_d}$ & $m_d=\frac{1}{w_d^2}$ & $W_0$ & $\frac{Q_0}{W_0}$ & $\delta_w$ \\
\midrule
I-1 & $(2\!+\!\sqrt3)+\frac{\sqrt{25\!+\!16\sqrt3}}{2}$ &
  $54.20$ & $0.0184$ &
  $18.92$ & $0.317$ & $9.85^\circ$ \\
  & $\dec{7.362}$ & & & & & \\[4pt]
I-2 & $(2\!+\!\sqrt3)-\frac{\sqrt{25\!+\!16\sqrt3}}{2}$ &
  $0.0104$ & $96.36$ &
  $0.861$ & $6.97$ & $6.68^\circ$ \\
  & $\dec{0.102}$ & & & & & \\[4pt]
II-1 & $\displaystyle -\frac{3+2\sqrt3}{6}$ &
  $\displaystyle\frac{7+4\sqrt3}{12}$ &
  $12(2\!-\!\sqrt3)^2$ &
  $\displaystyle\frac{2+\sqrt3}{3}$ &
  $18(2\!-\!\sqrt3)$ &
  $\mathbf{15.00}^\circ$ \\
  & $\dec{-1.077}$ & $\dec{1.161}$ & $\dec{0.862}$ & $\dec{1.244}$ & $\dec{4.82}$ & \\[4pt]
II-2 & $\displaystyle -\frac{21+2\sqrt3}{6}$ &
  $\displaystyle\frac{151+28\sqrt3}{12}$ &
  $\displaystyle\frac{12(151\!-\!28\sqrt3)}{143^2}$ &
  $\displaystyle\frac{14+3\sqrt3}{3}$ &
  $\displaystyle\frac{18(14\!-\!3\sqrt3)}{169}$ &
  $34.79^\circ$ \\
  & $\dec{-4.077}$ & $\dec{16.62}$ & $\dec{0.060}$ & $\dec{6.399}$ & $\dec{0.94}$ & \\[2pt]
\bottomrule
\end{tabular}
\caption{All Koide solutions for $K(1/m_d,\,1/m_s,\,1/m_b)$ at $\delta=15^\circ$.
  Cases~I use $(+w_s,+w_b)$, Cases~II use $(-w_s,+w_b)$.
  I-1 is the large-$|w_d|$ (small-$m_d$) branch; I-2 the small-$|w_d|$
  (large-$m_d$) branch.
  Grey numbers are decimal evaluations of the exact expressions.
  Only solution~II-1 has an angle in exact degrees.
  Units: $\sqrt{M_0}=1$.}
\label{tab:inverse-solutions}
\end{table}

\noindent
Solution~II-1 is distinguished: its discriminant is a perfect square,
its $m_d$ is a clean radical $12(2-\sqrt3)^2$, and its Koide angle
is $15^\circ$---the same angle as Tuple~1.  This self-referential
property singles it out among the four solutions.

\subsection{Galois duality: seed $(-1/q_s,\,+1/q_b)$ and its conjugate partner}

The alternative seed pair $(+1/q_c,\,+1/q_b)$ is, up to an overall sign, the
Galois conjugate
of $(-w_s,\,w_b)$ under $\sqrt3 \to -\sqrt3$ (which exchanges
$q_s\leftrightarrow q_c$). Since $q_s q_c = 3$ and $q_b = 2\sqrt3$,
the discriminant is again $\Delta = 9$ and the conjugate solutions are:
\[
  m_d^{(\pm)} = 12(7 \pm 4\sqrt3),
  \qquad
  m_d^{(\pm)} = \frac{12(151\pm 28\sqrt3)}{143^2}.
\]
Each solution is the Galois conjugate of the corresponding one from
Table~\ref{tab:inverse-solutions}, with the product
$m_d \cdot m_d^{(\text{conj})} = 144$
(resp.\ $144/143^2$).

The solution $m_d = 12(7+4\sqrt3)$ has Koide angle $\delta_w = 45^\circ$
---echoing Tuple~2, just as our $m_d = 12(7-4\sqrt3)$ echoes Tuple~1.
The solution $m_d = 12(151+28\sqrt3)/143^2$ has angle
$\pi - \arctan(4\sqrt3-5)\approx 117.4^\circ$ in a
$k=0,1,2$ input-ordered convention; in the Zenczykowski (magnitude-sorted)
convention this becomes $\approx 2.6^\circ$.  This angle is algebraically
exact in $\mathbb{Q}(\sqrt3)$ but is \emph{not} a rational multiple of
$\pi$.

In total, both seed pairs produce exactly three angles that are multiples
of $\pi/12$: the two forward-chain angles ($15^\circ$, $45^\circ$) plus
one seesaw angle that echoes one of them.  The apparent abundance of
``exact'' angles in the $k$-ordered convention is partly an artifact:
the phase depends on the $k$-assignment of the charges, and different
assignments give different (but equally algebraic) values.

\subsubsection{A complement near the lepton angle}\label{sec:galois-complement}

The Zenczykowski angle of $m_d = 12(151+28\sqrt3)/143^2$
(the Galois conjugate of our~II-2) is
\[
  \delta_{d+} = \arctan\frac{7\sqrt3-10}{47} \approx 2.59^\circ.
\]
Its complement with respect to $15^\circ$ has an exact tangent:
\[
  \tan(15^\circ - \delta_{d+}) = \frac{5\sqrt3-8}{3}\approx 0.2201,
  \qquad 15^\circ - \delta_{d+} = 12.41^\circ.
\]
This is close to the experimental lepton Koide angle
$\delta_{e\mu\tau} = 12.73^\circ$
($\tan\delta_{e\mu\tau}\approx 0.2260$), suggesting the decomposition
\[
  \delta_{d+} + \delta_{e\mu\tau} \approx 15^\circ.
\]
At the exact $\delta=15^\circ$ chain, the match is only approximate
(deficit $0.32^\circ$); along the continuation-preserving branch with
$\varepsilon < 0$, the gap
could close.

\section{Perturbation of the inverse-mass tuples}

When $\delta = 15^\circ + \varepsilon$ with $\varepsilon < 0$,
the double root $q_b = 2\sqrt{3}$ splits into two branches
$q_{b\pm} = 2\sqrt{3} \pm 6\sqrt{|\varepsilon|} + O(\varepsilon)$.
Each branch provides its own $w_{b\pm} = 1/q_{b\pm}$, and the
inverse-mass Koide analysis of \S5 must be repeated for each.
This doubles the number of solutions from 4 to 8.

\subsection{Exact projective II-1 identity along the deformed chain}
\label{sec:projective-ii}

Let $q_b$ denote either real continuation root of the direct Koide triple
$(-q_s,q_c,q_b)$.  Its Koide condition is
\[
  q_b^2 - 4(q_c-q_s)\,q_b + (q_s^2 + 4q_s q_c + q_c^2) = 0.
\]

The key observation is the projective identity
\[
  \left(-\frac{1}{q_s},\,\frac{1}{q_b},\,-\frac{q_c}{q_s q_b}\right)
  = -\frac{1}{q_s q_b}\,(q_b,\,-q_s,\,q_c).
\]
Since the Koide condition is homogeneous of degree two and symmetric under
permutations of the three entries, any overall rescaling or permutation of
a Koide triple is again Koide.  Therefore, once $(-q_s,q_c,q_b)$ is Koide,
\[
  \left(-\frac{1}{q_s},\,\frac{1}{q_b},\,-\frac{q_c}{q_s q_b}\right)
\]
is automatically Koide as well.

For completeness, the same statement can be checked directly at the level of
the inverse quadratic.  Consider the opposite-sign inverse seed
$(-1/q_s,\,+1/q_b,\,w_d)$.  Its Koide polynomial is
\[
  w_d^2 + 4\!\left(\frac{1}{q_s}-\frac{1}{q_b}\right) w_d
  + \left(\frac{1}{q_s^2} + \frac{4}{q_s q_b} + \frac{1}{q_b^2}\right)=0.
\]
Substituting
\[
  w_d = -\frac{q_c}{q_s q_b}
\]
\[
  \frac{q_b^2 - 4(q_c-q_s)\,q_b + q_s^2 + 4q_s q_c + q_c^2}{q_s^2 q_b^2}=0,
\]
whose numerator is exactly the direct polynomial above.
This direct substitution is redundant once the projective identity is
noticed, but it gives a useful quadratic cross-check.  The coincidence is
therefore a projective/factorised identity, not an identity of the centred
additive variables $z_0,z_i,w_0,\omega_i$.

Among the six permutations of $(-q_s,q_c,q_b)$, this is the only generic
direct/inverse proportionality with fixed second inverse entry $+1/q_b$:
the others collapse only on special loci such as $q_b=q_c$ or
$q_b^2=q_s q_c$.  By continuity from the exact $\varepsilon=0$ point, the
matched opposite-sign inverse root is precisely the II-1 branch on both
$b_+$ and $b_-$ continuations.

One isolated special locus is worth recording because it fixes an angle.
On the lower direct branch, imposing $q_{b-}=q_c$ in the direct quadratic
gives
\[
  -2q_c^2 + 8q_s q_c + q_s^2 = 0
  \qquad\Longrightarrow\qquad
  \frac{q_c}{q_s} = 2 + \frac{3\sqrt2}{2}.
\]
Numerically this occurs at
\[
  \varepsilon_{\rm eq} \approx -0.04158609,
  \qquad
  \delta_{\rm eq} \approx 12.61729^\circ.
\]
This does not create a second generic direct/inverse coincidence family, but
it marks a distinguished angle at which two direct entries become equal and
the existing II-1 permutation becomes partially degenerate.

\subsection{Numerical example: $\varepsilon = -0.001$}

At $\varepsilon = -0.001$ (i.e.\ $\delta \approx 14.94^\circ$),
the perturbed Tuple~1 charges are
$q_1 = 0.001$, $q_2 = 0.6326$, $q_3 = 2.3664$,
and the $b$ charge splits into $q_{b+} = 3.6573$ and $q_{b-} = 3.2778$.

\begin{table}[h]
\centering
\small
\begin{tabular}{@{}l l r r r r r@{}}
\toprule
Case & Branch & $w_d$ & $w_d^2$ & $m_d = 1/w_d^2$ & $W_0$ & $\delta_w$ \\
\midrule
I-1  & $b_+$ & $+7.301$ & $53.30$ & $0.019$ & $18.63$ & $10.07^\circ$ \\
I-1  & $b_-$ & $+7.454$ & $55.57$ & $0.018$ & $19.39$ & $9.63^\circ$ \\[3pt]
I-2  & $b_+$ & $+0.116$ & $0.013$ & $74.71$ & $0.862$ & $5.63^\circ$ \\
I-2  & $b_-$ & $+0.089$ & $0.008$ & $126.49$ & $0.867$ & $7.70^\circ$ \\[3pt]
II-1 & $b_+$ & $-1.023$ & $1.046$ & $0.956$ & $1.207$ & $17.05^\circ$ \\
II-1 & $b_-$ & $-1.141$ & $1.302$ & $0.768$ & $1.298$ & $12.87^\circ$ \\[3pt]
II-2 & $b_+$ & $-4.207$ & $17.70$ & $0.057$ & $6.756$ & $35.68^\circ$ \\
II-2 & $b_-$ & $-3.961$ & $15.69$ & $0.064$ & $6.095$ & $33.83^\circ$ \\
\bottomrule
\end{tabular}
\caption{All eight inverse-mass Koide solutions at $\varepsilon=-0.001$.
  Each solution from Table~\ref{tab:inverse-solutions} forks into a $b_+$
  and $b_-$ branch.}
\label{tab:inverse-perturbed}
\end{table}

\subsection{Observations}

\begin{itemize}
  \item Each of the four $\varepsilon=0$ solutions splits smoothly into
    a pair, one from each $q_b$ branch.
  \item The Cases~I discriminant $\Delta_{\rm I} = 25+16\sqrt3 + O(\sqrt{|\varepsilon|})$
    remains positive; Cases~II discriminant perturbs from $9$ but stays positive.
  \item The distinguished solution II-1 ($\delta_w = 15^\circ$ at $\varepsilon=0$)
    splits into $\delta_w = 17.05^\circ$ ($b_+$) and $\delta_w = 12.87^\circ$ ($b_-$).
    The $b_-$ branch is notable: at the experimental lepton angle
    $\delta_{e\mu\tau}\approx 12.73^\circ$
    ($\varepsilon\approx -0.040$), this branch would approach
    a value near the physical lepton angle itself.
\end{itemize}

\section{Comparison with experimental masses}

We now confront the algebraic results of \S\S1--6 with
the experimental quark and lepton masses.

\subsection{Reference angles for the standard Koide chain}

Using PDG~2024 $\overline{\rm MS}$ masses ($m_d = 4.67$~MeV,
$m_s = 93.4$~MeV, $m_c = 1.27$~GeV, $m_b = 4.18$~GeV
at $\mu=2$~GeV or $m_q(m_q)$; $m_t = 172.69$~GeV pole;
$m_e,m_\mu,m_\tau$ pole):

\begin{table}[h]
\centering
\small
\begin{tabular}{@{}l l r r l@{}}
\toprule
Tuple & Signs & $\delta$ & $Q$ & Note \\
\midrule
\multicolumn{5}{@{}l}{\emph{Standard Koide chain}} \\
$e,\mu,\tau$ (pole) & $(+,+,+)$ & $12.73^\circ$ & $1.50001$ & exact Koide \\
$(u,s,c)$ & $(+,+,+)$ & $13.28^\circ$ & $1.602$ & \\
$(0,s,c)$ [$m_u\!\to\!0$] & $(+,+,+)$ & $15.20^\circ$ & $1.505$ & near Koide \\
$(s,c,b)$ & $(-,+,+)$ & $37.21^\circ$ & $1.482$ & $\approx 3\delta_{e\mu\tau}$ \\
$(c,b,t)$ & $(+,+,+)$ & $3.93^\circ$ & $1.494$ & near Koide \\
\midrule
\multicolumn{5}{@{}l}{\emph{Other quark triples}} \\
$(d,s,b)$ & $(+,+,+)$ & $6.31^\circ$ & $1.367$ & \\
$(d,s,c)$ & $(+,+,+)$ & $12.33^\circ$ & $1.647$ & $\approx\delta_{e\mu\tau}$ \\
$(u,c,t)$ & $(+,+,+)$ & $4.26^\circ$ & $1.178$ & \\
\midrule
\multicolumn{5}{@{}l}{\emph{Inverse-mass tuples}} \\
$K(1/d,1/s,1/b)$ & $(+,+,+)$ & $10.70^\circ$ & $1.503$ & near Koide \\
$K(1/e,1/\mu,1/\tau)$ & $(+,+,+)$ & $2.73^\circ$ & $1.174$ & $\approx\delta_{d+}$ \\
$K(1/u,1/c,1/t)$ & $(+,+,+)$ & $1.91^\circ$ & $1.090$ & \\
\bottomrule
\end{tabular}
\caption{Koide angles and $Q$-parameters for chain tuples, other quark
  triples, and inverse-mass tuples (PDG~2024, $m_q(m_q)$ convention).
  Only $K(1/d,1/s,1/b)$ among inverse-mass tuples approaches $Q=3/2$.}
\label{tab:reference-angles}
\end{table}

\subsubsection*{Observations on Table~\ref{tab:reference-angles}}

\begin{itemize}
  \item The $(0,s,c)$ tuple at $\delta = 15.20^\circ$ confirms that
    the idealised $\delta = 15^\circ$ of \S1 is a good approximation.
    The $(-s,c,b)$ angle of $37.21^\circ$
    is close to $3\times 12.73^\circ = 38.2^\circ$, consistent with
    the factor-of-three relation $\delta_{scb}\approx 3\,\delta_{e\mu\tau}$.

  \item Among all inverse-mass tuples, \textbf{only $K(1/d,1/s,1/b)$
    approaches Koide}: $Q = 1.503$ ($0.2\%$ from $3/2$).
    The inverse lepton tuple $K(1/e,1/\mu,1/\tau)$ gives $Q = 1.17$
    and the inverse up-type $K(1/u,1/c,1/t)$ gives $Q=1.09$---both
    far from $3/2$.  This singles out the down-type quarks.

  \item The inverse lepton angle $\delta(K(1/e,1/\mu,1/\tau)) = 2.73^\circ$
    is remarkably close to the Galois conjugate seesaw angle
    $\delta_{d+} = 2.59^\circ$ from \S\ref{sec:galois-complement}.  While the inverse lepton
    tuple does \emph{not} satisfy Koide ($Q=1.17$), the near-coincidence
    of angles is intriguing.

  \item The ratios $\delta_w/\delta_{e\mu\tau} = 10.70/12.73 = 0.840$
    and $\delta_{e\mu\tau}/\delta_{0sc} = 12.73/15.20 = 0.837$
    are nearly identical, suggesting a geometric progression
    $\delta_w \approx \delta_{e\mu\tau}^2/\delta_{0sc}$.

  \item The $(d,s,c)$ tuple has $\delta = 12.33^\circ$, close to
    $\delta_{e\mu\tau} = 12.73^\circ$.  This is natural: in the exact
    chain, the $(0,s,c)$ tuple has $\delta = 15^\circ$, and
    replacing the zero by $m_d > 0$ pushes the angle downward.
    The $(d,s,c)$ tuple is the ``wrong light quark'' version of $(u,s,c)$
    ($\delta = 13.28^\circ$); both sit below $15^\circ$ and bracket
    the lepton angle.  The proximity of $(d,s,c)$ to $\delta_{e\mu\tau}$
    may reflect the same mechanism that links quark and lepton
    Koide angles in the chain.
\end{itemize}

\subsection{The inverse-mass tuple $K(1/m_d,\,1/m_s,\,1/m_b)$}

With all positive signs and PDG central values in the $m_q(m_q)$ convention:
\[
  Q = 1.503,\qquad \delta_w = 10.70^\circ.
\]
The deviation from Koide, $|Q - 3/2| = 3.3\times 10^{-3}$, is
larger than the lepton Koide ($1.4\times 10^{-5}$) but much
smaller than the direct $(d,s,b)$ tuple ($Q=1.37$).

To achieve $Q = 3/2$ exactly, one needs
$m_d = 4.61$~MeV, a $-1.4\%$ shift from the PDG central value
$4.67 \pm 0.48$~MeV---well within the $10\%$ experimental
uncertainty.

\subsection{RG running}

At one-loop QCD, all quark masses of the same charge run with a
common anomalous dimension, so their ratios---and hence the Koide
angle---are RG invariants within each flavour-number regime.
Threshold corrections at $\mu = m_b$ and $\mu = m_t$ introduce
small shifts.

\begin{table}[h]
\centering
\small
\begin{tabular}{@{}l r r r r@{}}
\toprule
Convention & $m_d$ (MeV) & $m_s$ (MeV) & $m_c$ (GeV) & $m_b$ (GeV) \\
\midrule
Common $\mu = 2$~GeV & $4.67$ & $93.4$ & $1.167$ & $4.667$ \\
$m_q(m_q)$ & $4.67$ & $93.4$ & $1.270$ & $4.180$ \\
\bottomrule
\end{tabular}

\medskip

\begin{tabular}{@{}l c r c r c r@{}}
\toprule
Convention
  & $\delta_{0sc}$ & $Q_{0sc}$
  & $\delta_{-scb}$ & $Q_{-scb}$
  & $\delta_w$ & $Q_w$
  \\
\midrule
Common $\mu = 2$~GeV
  & $15.93^\circ$ & $1.524$
  & $34.09^\circ$ & $1.453$
  & $10.79^\circ$ & $1.4992$
  \\
$m_q(m_q)$
  & $15.20^\circ$ & $1.505$
  & $37.21^\circ$ & $1.482$
  & $10.70^\circ$ & $1.5033$
  \\
\bottomrule
\end{tabular}
\caption{Koide angles for the chain tuples and the inverse-mass tuple
  $K(1/d,1/s,1/b)$ in two mass conventions.
  At 1-loop QCD the mass anomalous dimension is flavour-universal,
  so all angles are exactly scale-invariant within each convention---only
  the choice of convention matters, not the RG scale.
  The $\delta_{-scb}/\delta_{0sc}$ ratio is $2.14$ at common scale
  vs.\ $2.45$ in $m_q(m_q)$.}
\label{tab:rg-angles}
\end{table}

\subsubsection*{Observations}

\begin{itemize}
  \item \textbf{$(0,s,c)$ is above $15^\circ$ in both conventions.}
    $\delta_{0sc} = 15.93^\circ$ (common scale) or $15.20^\circ$
    ($m_q(m_q)$).  In both cases $\varepsilon > 0$ in the language
    of \S4, placing us past the bifurcation where the exact
    opposite-sign continuation has no real solutions.
    The experimental $(-s,c,b)$ tuple has $Q = 1.45$--$1.48$,
    close to but not equal to $3/2$---consistent with being
    near but past the critical angle.

  \item \textbf{1-loop QCD preserves all angles trivially.}
    The mass anomalous dimension is flavour-universal at one loop,
    so all quark mass ratios are exactly preserved under running.
    The angles in Table~\ref{tab:rg-angles} are strictly scale-invariant; only
    the choice of mass convention (common scale vs.\ $m_q(m_q)$)
    matters.  Yukawa contributions to the running, which are
    flavour-dependent (dominated by the top Yukawa), would break
    this invariance---particularly for $m_b$.  The question of how
    much $\delta_w$ and $\delta_{-scb}$ actually run under full
    SM RGEs is the natural next step for this programme.

  \item \textbf{$Q_w$ is closest to $3/2$ at common scale.}
    $Q_w = 1.4992$ vs.\ $1.5033$ in $m_q(m_q)$.
    The exact-Koide prediction at $\mu=2$~GeV
    is $m_d = 4.69$~MeV ($+0.4\%$ from PDG).
    With pole $m_b = 4.78$~GeV, $Q_w = 1.4983$.
\end{itemize}

\section{Which branch are we on?}

The exact chain at $\delta = 15^\circ$ has a double root for $m_b$
and four seesaw solutions for $m_d$.  At the physical point,
$\delta_{0sc}$ from Table~\ref{tab:rg-angles} lies slightly above
$15^\circ$, i.e.\ on the $\varepsilon>0$ side of the bifurcation.  For branch
selection, however, we follow the real continuation-preserving side
$\varepsilon < 0$, where the double root splits into $q_{b\pm}$ and
each seesaw solution forks into two branches, giving eight candidates.
We now ask: for which branches does the inverse-mass Koide angle
$\delta_w$ match the experimental value $\approx 10.70^\circ$,
and does the rest of the chain produce sensible mass ratios?

\subsection{Matching $\delta_w = 10.70^\circ$}

Here $\delta_w$ is the Zenczykowski angle of the \emph{mass triple}
$(m_d,m_s,m_b)$ generated by a given signed seed, with the sign pattern
recorded separately in the branch label.  The observable inverse-mass tuple
$K(1/d,1/s,1/b)$ uses all positive entries; the signed seeds are an auxiliary
device for selecting which exact Koide continuation produces that positive
mass spectrum.  The comparison is therefore between two descriptions of the
same masses, not between two incompatible observables.

A numerical scan over $\varepsilon \in (-0.15,\,0)$ finds exactly
three branches where $\delta_w$ passes through $10.70^\circ$:

\begin{table}[h]
\centering
\small
\begin{tabular}{@{}l r r r r r r r r@{}}
\toprule
Branch & $\varepsilon$ & $\delta_1$ & $m_d/m_s$ & $m_s/m_c$ & $m_b/m_c$
  & $\delta_{-scb}$ & $m_d$ & $m_b$ \\
  & (rad) & & & & & & (MeV) & (MeV) \\
\midrule
\textbf{I-1/$b_+$} & $\mathbf{-0.0137}$ & $\mathbf{14.2^\circ}$ & $\mathbf{0.050}$ & $\mathbf{0.067}$ & $\mathbf{3.16}$
  & $\mathbf{37.8^\circ}$ & $\mathbf{4.3}$ & $\mathbf{4033}$ \\
I-2/$b_-$ & $-0.0169$ & $14.0^\circ$ & $944$ & $0.066$ & $1.34$
  & $54.2^\circ$ & $80021$ & $1707$ \\
II-1/$b_-$ & $-0.1101$ & $8.7^\circ$ & $0.58$ & $0.041$ & $0.58$
  & $49.3^\circ$ & $31$ & $763$ \\
\bottomrule
\end{tabular}
\caption{Branches matching $\delta_w\approx 10.70^\circ$.
  Physical masses obtained by scaling $Z_0 = m_s + m_c$ to $1363$~MeV.
  Experimental ratios: $m_d/m_s\approx 0.050$, $m_b/m_c\approx 3.3$.}
\label{tab:branch-match}
\end{table}

\subsection{Filtering by mass ratios}

\begin{itemize}
  \item \textbf{I-2/$b_-$: excluded.}  $m_d/m_s = 944$: the down quark
    is heavier than the bottom.
  \item \textbf{II-1/$b_-$: excluded.}  $m_b/m_c = 0.58$: the bottom
    is lighter than charm.
  \item \textbf{I-1/$b_+$: viable.}  All mass ratios are close to
    experiment:
    \begin{itemize}
      \item $m_d/m_s = 0.050$ (exp: $0.050$),
      \item $m_s/m_c = 0.067$ (exp: $0.074$),
      \item $m_b/m_c = 3.16$ (exp: $3.29$),
      \item $\delta_{-scb} = 37.8^\circ$ (vs.\ $3\delta_{e\mu\tau} = 38.2^\circ$).
    \end{itemize}
\end{itemize}

\subsection{The surviving solution}

The unique viable branch is \textbf{I-1/$b_+$}: the same-sign seesaw
$(+w_s,\,+w_b,\,w_d)$ with the heavy $b$-branch, at
$\varepsilon = -0.0137$ ($\delta_1 = 14.2^\circ$).

At this point the chain predicts (scaling to $Z_0 = 1363$~MeV):
\[
  m_u = 0.04~\text{MeV},\quad
  m_d = 4.3~\text{MeV},\quad
  m_s = 86~\text{MeV},\quad
  m_c = 1277~\text{MeV},\quad
  m_b = 4033~\text{MeV}.
\]
All except the very small $m_u$ lie within about $1$--$2\sigma$ of PDG
central values; $m_u$ comes out parametrically too light in this fit.
The $(-s,c,b)$ angle $37.8^\circ$ is within $1\%$ of
$3\delta_{e\mu\tau} = 38.2^\circ$, preserving the Zenczykowski
relation.

Note that $\delta_1 = 14.2^\circ$ is the angle of the \emph{first}
tuple $(0,s,c)$ at this $\varepsilon$ --- intermediate between the
idealised $15^\circ$ and the experimental lepton angle $12.73^\circ$.
The seesaw lives on the same-sign, $b_+$ branch: the only combination
that simultaneously satisfies the inverse-mass Koide equation and
produces a realistic quark spectrum.

\subsection{Near-compatibility of two constraints}

The system has one free parameter ($\varepsilon$) and two
natural constraints:
\begin{enumerate}
  \item[(A)] $\delta_w = 10.70^\circ$ (experimental inverse-mass angle),
  \item[(B)] $\delta_{-scb} = 3\,\delta_{e\mu\tau} = 38.20^\circ$
    (Zenczykowski relation).
\end{enumerate}
Fixing (A) gives $\varepsilon_A = -0.0137$, $\delta_{-scb} = 37.79^\circ$
(miss $-0.41^\circ$).
Fixing (B) gives $\varepsilon_B = -0.0121$, $\delta_w = 10.65^\circ$
(miss $-0.05^\circ$).
The gap $|\varepsilon_A - \varepsilon_B| = 0.0016$~rad ($0.09^\circ$)
is small: the two constraints are nearly compatible, with a residual
of $\sim 1\%$ in $\delta_{-scb}$.
This residual could plausibly be absorbed by Yukawa corrections
to $m_b$ running, which we have not included.

\section{Other tuples}

\subsection{The chiral tuple $(m_d,\,0,\,m_s)$}

Descending the chain from $(0,\,m_s,\,m_c)$ with $m_u = 0$,
the previous link is $(m_d,\,0,\,m_s)$.  The Koide equation
gives $q_d = (2\pm\sqrt3)\,q_s$, i.e.
\[
  q_{d+} = (2+\sqrt3)\,q_s = \frac{3+\sqrt3}{2} = q_c,
  \qquad
  q_{d-} = (2-\sqrt3)\,q_s = \frac{9-5\sqrt3}{2}.
\]
The ``$+$'' solution gives back the charm charge---the chain
\emph{loops}: descending from $(0,s,c)$ recovers $m_c$ as one branch.
The physical solution is $q_{d-}$, with
\[
  m_{d-} = (7-4\sqrt3)\,m_s,
  \qquad
  M_0^{(d,0,s)} = (7-4\sqrt3)\,M_0^{(0,s,c)}.
\]
Both solutions have $\delta = 15^\circ$ exactly (a zero entry
always reproduces the seed angle).  The scale parameter
$M_0$ is multiplied by $(2-\sqrt3)^2 = 7-4\sqrt3$---the same
factor that relates consecutive Koide tuples in the chain.

\medskip\noindent
Replacing $m_d$ by the \emph{chiral mass} $m_u + m_d$
(the quantity entering the chiral condensate) gives a physical
near-Koide tuple:
\[
  (m_u\!+\!m_d,\; 0,\; m_s) = (6.83,\;0,\;93.4)\text{ MeV}:
  \qquad Q = 1.504,\quad \delta = 15.15^\circ.
\]
The chiral combination works better than $m_d$ alone
(which gives $Q=1.43$, $\delta=12.3^\circ$): it sits just
above $15^\circ$, consistent with the bifurcation analysis of \S4.

\subsection{The $(s,\,b,\,t)$ tuple}

The top charge from $(c,b,t)$ at $\delta = 15^\circ$ simplifies
via $\sqrt{4+2\sqrt3} = 1+\sqrt3$:
\[
  q_t = \frac{15+19\sqrt3}{2},
  \qquad
  m_t = \frac{654+285\sqrt3}{2}.
\]
The key mass ratios in the exact chain are:
\[
  \frac{m_t}{m_c} = 151-28\sqrt3,
  \qquad
  \frac{m_b}{m_c} = 16-8\sqrt3 = 8(2-\sqrt3),
  \qquad
  \frac{m_c}{m_s} = 7+4\sqrt3.
\]
Note that $151$ and $28$ are the same integers appearing
in the inverse-mass seesaw solution
$m_d^{(\text{II-2})} = 12(151-28\sqrt3)/143^2$.

The $(s,b,t)$ tuple gives
\[
  Q_{sbt} = 1.342,\qquad \delta_{sbt} = 6.38^\circ
  \qquad\text{(exact chain)},
\]
far from $3/2$. With PDG masses: $Q = 1.36$, $\delta = 7.2^\circ$.
Neither sign pattern helps ($Q_{-s,b,t} = 1.22$).
This tuple is genuinely non-Koide at every scale.

\subsection{The corrected inverse tuple $K(1/s,\,1/c,\,\pm 1/t)$}

The tuple $K(1/s,\,1/c,\,0)$ is the inversion of $(0,\,s,\,c)$
and has $\delta = 15^\circ$, $Q = 3/2$ exactly.
The inverse top charge is
\[
  \frac{1}{q_t} = \frac{19\sqrt3-15}{429},
  \qquad 429 = 3\times 11\times 13.
\]
Adding $\pm 1/q_t$ as a correction to the zero entry:

\begin{center}
\begin{tabular}{@{}l r r r r@{}}
\toprule
Tuple & $\delta$ (exact) & $Q$ (exact) & $\delta$ (PDG) & $Q$ (PDG) \\
\midrule
$K(1/s,\,1/c,\,0)$ & $15.00^\circ$ & $1.500$ & $15.20^\circ$ & $1.505$ \\
$K(1/s,\,1/c,\,+\!1/t)$ & $13.78^\circ$ & $1.562$ & $14.13^\circ$ & $1.560$ \\
$K(1/s,\,1/c,\,-\!1/t)$ & $16.17^\circ$ & $1.437$ & $16.23^\circ$ & $1.450$ \\
\bottomrule
\end{tabular}
\end{center}

\noindent
Adding $+1/q_t$ pushes $\delta$ \emph{below} $15^\circ$ (toward
the lepton angle); adding $-1/q_t$ pushes it above.
Neither achieves $Q=3/2$, but the $(+)$ correction brings $\delta$
closer to $\delta_{e\mu\tau}=12.73^\circ$ ($13.78^\circ$ at the
chain point, $14.13^\circ$ with PDG masses).
The correction is small ($\sim 1^\circ$) because
$1/q_t = 0.042 \ll 1/q_c = 0.423$---the top mass suppresses the
third entry by an order of magnitude.

\section{Meson Koide tuples}

The Koide equation can be applied to meson masses.
We scan all triples of pseudoscalar mesons with both sign
patterns $(+,+,+)$ and $(-,+,+)$.

\subsection{Best pseudoscalar triples}

\begin{table}[h]
\centering
\small
\begin{tabular}{@{}l l r r l@{}}
\toprule
Tuple & Signs & $\delta$ & $Q$ & $|Q\!-\!3/2|$ \\
\midrule
\multicolumn{5}{@{}l}{\emph{Type $(-\pi,\,D_s,\,X)$}} \\
$(\pi^0,\, D_s,\, \eta_c)$ & $(-,+,+)$ & $51.73^\circ$ & $1.5006$ & $6\times 10^{-4}$ \\
$(\pi^\pm,\, D_s,\, B^\pm)$ & $(-,+,+)$ & $40.79^\circ$ & $1.4984$ & $1.6\times 10^{-3}$ \\
$(\pi^\pm,\, D_s,\, B_s)$ & $(-,+,+)$ & $40.49^\circ$ & $1.4978$ & $2.2\times 10^{-3}$ \\
$(\pi^0,\, D_s,\, B_c)$ & $(-,+,+)$ & $37.65^\circ$ & $1.4962$ & $3.8\times 10^{-3}$ \\[3pt]
\multicolumn{5}{@{}l}{\emph{Type $(-\pi,\,D,\,B)$}} \\
$(\pi^0,\, D^\pm,\, B^\pm)$ & $(-,+,+)$ & $39.89^\circ$ & $1.4929$ & $7.1\times 10^{-3}$ \\
$(\pi^\pm,\, D^\pm,\, B^\pm)$ & $(-,+,+)$ & $39.94^\circ$ & $1.4864$ & $1.4\times 10^{-2}$ \\[3pt]
\multicolumn{5}{@{}l}{\emph{Other}} \\
$(\pi^\pm,\, \pi^0,\, \eta_b)$ & $(+,+,+)$ & $0.11^\circ$ & $1.4981$ & $1.9\times 10^{-3}$ \\
\bottomrule
\end{tabular}
\caption{Best pseudoscalar meson Koide tuples (PDG~2024 masses).
  The $(-\pi, D_s, \eta_c)$ tuple rivals the lepton Koide in
  precision.  All require a sign flip on the pion.}
\end{table}

\noindent
The $(-\pi^0,\, D_s,\, \eta_c)$ tuple is remarkable:
$|Q - 3/2| = 6\times 10^{-4}$, only $40\times$ worse than the
lepton Koide.  The pattern is consistent: \emph{all} good meson
triples involve a sign flip on the pion (the lightest member)
and pair it with a charmed and a heavier meson.
The $D_s$ (containing an $s$ quark) works better than the $D$,
suggesting the strange content improves the match.

\subsection{Vector mesons}

No triple of vector mesons ($\rho, \omega, K^*, \phi, D^*, B^*,
J/\psi, \Upsilon$) comes close to Koide: the best is
$(-\rho,\, B_s^*,\, \Upsilon)$ with $Q = 1.31$.
Vector mesons do not form Koide tuples.

\subsection{Inverse-mass mesons}

Among pure meson triples, no inverse-mass tuple approaches Koide:
the best is $K(1/\pi^0,\,1/B_c,\,1/\eta_b)$ with $Q = 1.55$.

\subsection{Including $W$, $Z$, and $H$}

Adding the electroweak bosons ($M_W = 80.36$, $M_Z = 91.19$,
$M_H = 125.2$~GeV) to the scan produces several near-Koide
mixed meson-boson tuples:

\begin{center}
\small
\begin{tabular}{@{}l l r r r@{}}
\toprule
Tuple & Signs & $\delta$ & $Q$ & $|Q\!-\!3/2|$ \\
\midrule
$K(1/\pi^0,\,1/\eta_c,\,1/W)$ & $(+,+,+)$ & $9.66^\circ$ & $1.5013$ & $1.3\times 10^{-3}$ \\
$(\pi^0,\,B_c,\,H)$ & $(+,+,+)$ & $10.75^\circ$ & $1.5024$ & $2.4\times 10^{-3}$ \\
$(D^\pm,\,D^0,\,H)$ & $(+,+,+)$ & $0.01^\circ$ & $1.5033$ & $3.3\times 10^{-3}$ \\
$(\eta',\,D^0,\,Z)$ & $(+,+,+)$ & $2.29^\circ$ & $1.5047$ & $4.7\times 10^{-3}$ \\
$(-B_c,\,Z,\,H)$ & $(-,+,+)$ & $53.70^\circ$ & $1.4931$ & $6.9\times 10^{-3}$ \\
\bottomrule
\end{tabular}
\end{center}

\noindent
The pure boson triple $(W,Z,H)$ gives $Q = 2.97$---not Koide.
But mixed meson-boson tuples can achieve $|Q-3/2| \sim 10^{-3}$,
comparable to $K(1/d,1/s,1/b)$.  The inverse-mass tuple
$K(1/\pi^0,1/\eta_c,1/W)$ is particularly clean: all positive signs,
$\delta = 9.7^\circ$, and $Q$ within $0.1\%$ of $3/2$.
Whether these reflect structure or coincidence across
the meson-boson mass landscape requires further investigation.

\subsection{Baryons: direct and inverse}

Using PDG~2025 baryon masses, we scan all triples among the
22 ground-state baryons ($p$, $n$, $\Lambda$, $\Sigma^{+,0,-}$,
$\Xi^{0,-}$, $\Omega^-$, $\Lambda_c^+$, $\Sigma_c^{++,+,0}$,
$\Xi_c^{+,0}$, $\Omega_c^0$, $\Lambda_b^0$, $\Sigma_b^{+,-}$,
$\Xi_b^{0,-}$, $\Omega_b^-$).

\medskip\noindent
\textbf{Direct:}  The best triple is
$(-p,\,\Sigma_b^-,\,\Omega_b^-)$ with $Q = 1.19$---far from $3/2$.
No baryon triple with any sign pattern approaches Koide.

\medskip\noindent
\textbf{Inverse:}  The best is
$K(1/p,\,1/\Sigma_b^-,\,1/\Omega_b^-)$ with $Q = 2.45$---even
farther.

\medskip\noindent
The Koide property appears specific to: charged leptons (exact),
the quark chain $\{(0sc),(-scb),(cbt)\}$ and its inverse-mass
descendant $K(1/d,1/s,1/b)$ (approximate), and certain
pseudoscalar meson triples with sign-flipped pion (approximate).
Baryons do not participate.

\section{Preon charge patterns}

For any signed direct tuple with charges
$q_i = s_i\sqrt{m_i}$, one can extract a common piece $z_0$ and a traceless
part $z_i$ by
\[
  q_i = z_0 + z_i,
  \qquad
  z_0 = \frac{q_1+q_2+q_3}{3},
  \qquad
  z_1+z_2+z_3 = 0.
\]
Likewise, for any inverse tuple with signed inverse charges
$w_i = s_i/\sqrt{m_i}$,
\[
  w_i = w_0 + \omega_i,
  \qquad
  w_0 = \frac{w_1+w_2+w_3}{3},
  \qquad
  \omega_1+\omega_2+\omega_3 = 0.
\]
The observed masses are recovered as
\[
  m_i = (z_0 + z_i)^2,
  \qquad
  \frac{1}{m_i} = (w_0 + \omega_i)^2.
\]
If one wishes to interpret the constituents themselves as ``mean preon
masses'', the corresponding scales are obtained by squaring the direct charges
or inverse-squaring the inverse charges:
$z_0^2,z_i^2$ and $1/w_0^2,1/\omega_i^2$.

It is also convenient to define the centred ratios
\[
  R_z \equiv \frac{z_1^2+z_2^2+z_3^2}{3z_0^2},
  \qquad
  R_w \equiv \frac{\omega_1^2+\omega_2^2+\omega_3^2}{3w_0^2}.
\]
Since $Q = 9z_0^2/(3z_0^2+z_1^2+z_2^2+z_3^2)$ (and similarly for $w$),
one has
\[
  R = \frac{3}{Q} - 1.
\]
Thus exact Koide corresponds simply to $R=1$, while near-Koide tuples have
$R$ close to $1$.

\subsection{Direct tuples}

For the two exact quark-chain anchors, the decomposition is especially simple:
\[
  (0,s,c):
  \qquad
  z_0 = 1,
  \qquad
  (z_1,z_2,z_3)
  = \left(-1,\frac{1-\sqrt3}{2},\frac{1+\sqrt3}{2}\right),
\]
\[
  (-s,c,b):
  \qquad
  z_0 = \sqrt3,
  \qquad
  (z_1,z_2,z_3)
  = \left(-\frac{3+\sqrt3}{2},\frac{3-\sqrt3}{2},\sqrt3\right).
\]
It is also useful to record an alternative \emph{factorised} normalization
for the exact tuples.  Writing $q_i = \bar z\,\hat z_i$ with
$\bar z = (q_1+q_2+q_3)/3$ and $\hat z_1+\hat z_2+\hat z_3=3$, and defining
\[
  a \equiv \frac{3-\sqrt3}{2},
  \qquad
  b \equiv \frac{3+\sqrt3}{2},
\]
one has the exact reduced shapes
\[
  (0,s,c) = 1\cdot(0,a,b),\qquad
  (d,0,s)_{\rm phys} = (2-\sqrt3)\cdot(a,0,b),\qquad
  (-s,c,b) = \sqrt3\cdot(a-1,b-1,2).
\]
This factorised notation is not the centred preon decomposition used in the
tables below, but it makes the exact shape relations transparent.

The same centred decomposition exists for the other exact and empirical tuples:

\begin{table}[h]
\centering
\tiny
\setlength{\tabcolsep}{2pt}
\begin{tabular}{@{}l c c c c c@{}}
\toprule
Tuple & Signs & preon masses & $z_0$ & $(z_1,z_2,z_3)$ & $R_z$ \\
\midrule
\multicolumn{6}{@{}l}{\emph{Exact tuples, $\sqrt{M_0}=1$ so masses are unitless}} \\
exact $(0,s,c)$ & $(+,+,+)$ & $(1.000;1.000,0.134,1.866)$ & $1.000$ & $(-1.000,-0.366,1.366)$ & $1.000$ \\
exact $(-s,c,b)$ & $(-,+,+)$ & $(3.000;5.598,0.402,3.000)$ & $1.732$ & $(-2.366,0.634,1.732)$ & $1.000$ \\
exact $(d,0,s)$ [phys.] & $(+,+,+)$ & $(0.072;0.010,0.072,0.134)$ & $0.268$ & $(-0.098,-0.268,0.366)$ & $1.000$ \\
exact $(c,b,t)$ & $(+,+,+)$ & $(98.57;57.19,41.78,196.7)$ & $9.928$ & $(-7.562,-6.464,14.026)$ & $1.000$ \\
exact $(s,b,t)$ & $(+,+,+)$ & $(87.44;75.98,34.65,213.3)$ & $9.351$ & $(-8.717,-5.887,14.604)$ & $1.235$ \\
\midrule
\multicolumn{6}{@{}l}{\emph{Empirical / fitted tuples, masses in MeV}} \\
fit $(-s,c,b)$ [I-1/$b_+$] & $(-,+,+)$ & $(899.3;1541.6,33.02,1123.4)$ & $29.989$ & $(-39.263,5.746,33.517)$ & $1.000$ \\
chiral $(u+d,0,s)$ & $(+,+,+)$ & $(16.75;2.188,16.75,31.05)$ & $4.093$ & $(-1.479,-4.093,5.572)$ & $0.995$ \\
$(e,\mu,\tau)$ & $(+,+,+)$ & $(313.8;289.0,55.30,597.2)$ & $17.716$ & $(-17.001,-7.437,24.437)$ & $1.000$ \\
$(u,s,c)$ & $(+,+,+)$ & $(243.1;199.4,35.12,401.9)$ & $15.590$ & $(-14.121,-5.926,20.047)$ & $0.873$ \\
$(0,s,c)$ & $(+,+,+)$ & $(228.0;228.0,29.55,421.8)$ & $15.100$ & $(-15.100,-5.436,20.537)$ & $0.993$ \\
$(-s,c,b)$ & $(-,+,+)$ & $(912.6;1589.8,29.47,1186.4)$ & $30.209$ & $(-39.873,5.429,34.444)$ & $1.025$ \\
$(d,s,c)$ & $(+,+,+)$ & $(250.3;186.6,37.90,392.7)$ & $15.821$ & $(-13.660,-6.156,19.816)$ & $0.822$ \\
$(d,s,b)$ & $(+,+,+)$ & $(649.9;544.4,250.5,1533.5)$ & $25.493$ & $(-23.332,-15.828,39.160)$ & $1.194$ \\
$(c,b,t)$ & $(+,+,+)$ & $(29566.8;18581.2,11512.7,59345.8)$ & $171.9$ & $(-136.3,-107.3,243.6)$ & $1.008$ \\
$(u,c,t)$ & $(+,+,+)$ & $(22767.5;22326.1,13283.0,70050.7)$ & $150.9$ & $(-149.4,-115.3,264.7)$ & $1.547$ \\
PDG $(s,b,t)$ & $(+,+,+)$ & $(26664.4;23601.6,9729.7,63638.9)$ & $163.3$ & $(-153.6,-98.639,252.3)$ & $1.212$ \\
\bottomrule
\end{tabular}
\caption{Direct charge decompositions $q_i = z_0 + z_i$ for the quark/lepton
  tuples discussed in the text.  The mass column lists the preon component masses
  $(z_0^2; z_1^2,z_2^2,z_3^2)$.  Exact Koide corresponds to $R_z=1$.
  Empirical rows use the same PDG conventions as
  Table~\ref{tab:reference-angles}.}
\end{table}

The direct patterns suggest a clear hierarchy.  The lepton tuple and the
exact/fitted quark-chain tuples all sit very close to $R_z=1$.
The empirical $(0,s,c)$ and $(-s,c,b)$ tuples remain near this surface,
whereas $(d,s,b)$, $(u,c,t)$ and $(s,b,t)$ drift farther away.
A simple scan over the centred charges and their preon masses shows that the
quark chain is unusually rigid.  In the exact solutions, passing from
$(0,s,c)$ to $(-s,c,b)$ multiplies each centred preon mass by a factor of
$3$.  Empirically the pattern softens but survives in recognisable form:
the common piece nearly doubles, $z_0:15.10\to 30.21$, so its preon mass
nearly quadruples, $z_0^2:228\to 913$ MeV.  Few other families exhibit such
a clean component-wise scaling.

\subsection{Inverse tuples}

For inverse tuples there are two natural conventions.  One may either
decompose the observable positive inverse-mass triple
$K(1/m_1,1/m_2,1/m_3)$, or keep the signed inverse charges that define a
specific seesaw branch.  To preserve the branch information, the exact
I/II solutions below are written with their signed inverse charges in the
$(d,s,b)$ order.

\begin{table}[h]
\centering
\tiny
\setlength{\tabcolsep}{2pt}
\begin{tabular}{@{}l c c c c c@{}}
\toprule
Tuple / seed & Signs & preon masses & $w_0$ & $(\omega_1,\omega_2,\omega_3)$ & $R_w$ \\
\midrule
\multicolumn{6}{@{}l}{\emph{Exact inverse solutions, $\sqrt{M_0}=1$ so preon masses are unitless}} \\
I-1 & $(+,+,+)$ & $(0.106;0.054,0.445,0.129)$ & $3.076$ & $(4.286,-1.499,-2.787)$ & $1.000$ \\
I-2 & $(+,+,+)$ & $(2.324;3.257,1.178,7.413)$ & $0.656$ & $(-0.554,0.921,-0.367)$ & $1.000$ \\
II-1 & $(-,-,+)$ & $(1.608;12.000,1.608,0.862)$ & $-0.789$ & $(-0.289,-0.789,1.077)$ & $1.000$ \\
II-2 & $(-,-,+)$ & $(0.313;0.191,22.392,0.232)$ & $-1.789$ & $(-2.289,0.211,2.077)$ & $1.000$ \\
exact $K(1/s,1/c,0)$ & $(+,+,0)$ & $(2.250;1.206,16.794,2.250)$ & $0.667$ & $(0.911,-0.244,-0.667)$ & $1.000$ \\
exact $K(1/s,1/c,+1/t)$ & $(+,+,+)$ & $(2.159;1.244,15.031,2.450)$ & $0.681$ & $(0.897,-0.258,-0.639)$ & $0.920$ \\
exact $K(1/s,1/c,-1/t)$ & $(+,+,-)$ & $(2.347;1.170,18.887,2.073)$ & $0.653$ & $(0.925,-0.230,-0.694)$ & $1.088$ \\
\midrule
\multicolumn{6}{@{}l}{\emph{Empirical / fitted inverse tuples, preon masses in MeV}} \\
PDG $K(1/d,1/s,1/b)$ & $(+,+,+)$ & $(26.60;13.84,122.3,31.41)$ & $0.194$ & $(0.269,-0.090,-0.178)$ & $0.996$ \\
PDG $K(1/e,1/\mu,1/\tau)$ & $(+,+,+)$ & $(3.896;1.256,5.968,4.288)$ & $0.507$ & $(0.892,-0.409,-0.483)$ & $1.554$ \\
PDG $K(1/u,1/c,1/t)$ & $(+,+,+)$ & $(17.81;5.085,22.915,18.177)$ & $0.237$ & $(0.443,-0.209,-0.235)$ & $1.753$ \\
PDG $K(1/s,1/c,0)$ & $(+,+,0)$ & $(520.2;281.3,4014.0,520.2)$ & $0.044$ & $(0.060,-0.016,-0.044)$ & $0.993$ \\
PDG $K(1/s,1/c,+1/t)$ & $(+,+,+)$ & $(501.7;289.0,3635.1,560.5)$ & $0.045$ & $(0.059,-0.017,-0.042)$ & $0.923$ \\
PDG $K(1/s,1/c,-1/t)$ & $(+,+,-)$ & $(539.8;273.8,4455.3,484.1)$ & $0.043$ & $(0.060,-0.015,-0.045)$ & $1.069$ \\
fit $K(1/d,1/s,1/b)$ [I-1/$b_+$] & $(+,+,+)$ & $(24.52;12.73,112.9,28.85)$ & $0.202$ & $(0.280,-0.094,-0.186)$ & $0.998$ \\
\bottomrule
\end{tabular}
\caption{Inverse charge decompositions $w_i = w_0 + \omega_i$.
  The mass column lists the inverse preon-component masses
  $(1/w_0^2; 1/\omega_1^2,1/\omega_2^2,1/\omega_3^2)$.
  The exact branch solutions use the signed inverse charges that generate
  the seesaw branches; the empirical inverse tuples use the observable
  positive inverse masses.}
\end{table}

The distinguished inverse down-type solutions again sit near $R_w=1$:
exact II-1 gives $R_w=1$ exactly, the empirical
$K(1/d,1/s,1/b)$ gives $R_w=0.996$, and the fitted surviving branch gives
$R_w=0.998$.  By contrast, the inverse lepton and inverse up-type tuples
move far away from this centred surface.
In the corresponding factorised notation $w_i=\bar w\,\hat w_i$, the stronger
statement proved in \S\ref{sec:projective-ii} is that the continued II-1 branch is projectively
identical to the direct continued branch for both $q_{b\pm}$:
\[
  \left(-\frac{1}{q_s},\,\frac{1}{q_b},\,-\frac{q_c}{q_s q_b}\right)
  = -\frac{1}{q_s q_b}\,(q_b,\,-q_s,\,q_c).
\]
At the exact $\delta=15^\circ$ point this reduces to the simple radical shape
\[
  ( -w_s,\,+w_b,\,w_d )_{\rm II-1}
  = -\frac{3+\sqrt3}{6}\cdot\left(2,\frac{1-\sqrt3}{2},\frac{1+\sqrt3}{2}\right),
\]
which is the permutation partner of the exact direct $(-s,c,b)$ shape.
This coincidence lives at the factorised/projective level; it is not an exact
identity of the centred additive preon components.
Beyond this exact one-parameter II-1 family, a continuity-tracked scan of the
local $\varepsilon<0$ branches
reveals one additional off-diagonal recurrence: on the $I$-2/$b_-$ branch,
at
\[
  \varepsilon_\star \approx -0.01014665808,
\]
the factorised reduced shape coincides with the exact $I$-1 shape,
again up to permutation.  Unlike the weaker drifts seen near the edge of a
finite scan window, this is an interior recurrence and therefore appears to
reflect a genuine algebraic relation rather than a boundary artefact.
The same scan finds many exact sum and difference identities among
$K(1/s,1/c,0)$ and $K(1/s,1/c,\pm 1/t)$, but these largely reflect the
fact that the centred inverse charges differ only by the small top entry.
This inverse family is therefore structured, but not in a mysterious way:
the algebra is real, yet much of it is inherited from the simple
$0,\pm 1/t$ deformation of the same seed.

\subsection{Pion-led composite patterns}

The retained meson tuples of \S11 can also be centred in the same way.
If one interprets the common piece $z_0$ as a string-like background
charge shared by the tuple, the pion-led direct tuples exhibit a fairly
stable background scale together with a large negative pion residual:

\begin{table}[h]
\centering
\scriptsize
\setlength{\tabcolsep}{3pt}
\begin{tabular}{@{}l c c c c c@{}}
\toprule
Tuple & Type & preon masses [MeV] & $Q$ & common piece & residuals \\
\midrule
$(0,\pi^0,D^\pm)$ & seed & $(334.4;334.4,44.46,622.7)$ & $1.5012$ & $z_0 = 18.286$ & $(-18.286,-6.668,24.954)$ \\
$(0,\pi^\pm,D_s)$ & seed & $(350.7;350.7,47.79,657.4)$ & $1.4973$ & $z_0 = 18.727$ & $(-18.727,-6.913,25.639)$ \\
$(-\pi^0,D_s,\eta_c)$ & direct & $(848.2;1659.9,232.3,650.3)$ & $1.5006$ & $z_0 = 29.124$ & $(-40.742,15.242,25.501)$ \\
$(-\pi^\pm,D_s,B^\pm)$ & direct & $(1229.9;2198.2,86.41,1412.9)$ & $1.4984$ & $z_0 = 35.070$ & $(-46.884,9.296,37.589)$ \\
$(-\pi^\pm,D_s,B_s)$ & direct & $(1244.0;2216.9,82.73,1443.1)$ & $1.4978$ & $z_0 = 35.270$ & $(-47.084,9.096,37.989)$ \\
$(-\pi^0,D_s,B_c)$ & direct & $(1392.9;2395.1,49.61,1755.3)$ & $1.4962$ & $z_0 = 37.322$ & $(-48.940,7.044,41.896)$ \\
$(\pi^\pm,\pi^0,\eta_b)$ & direct & $(1610.1;801.6,812.7,3228.5)$ & $1.4981$ & $z_0 = 40.126$ & $(-28.312,-28.508,56.820)$ \\
$K(1/\pi^0,1/\eta_c,1/W)$ & inverse & $(772.9;398.3,3205.5,950.1)$ & $1.5013$ & $w_0 = 0.036$ & $(0.050,-0.018,-0.032)$ \\
\bottomrule
\end{tabular}
\caption{Centred charge decompositions for the zero-entry pion seeds and the
  retained pion-led composite tuples.  The seed rows isolate the
  $(0,\pi,D)$-type background before the sign-flipped heavy continuation to
  $(-\pi,D,X)$.  The mass column lists the preon masses of the common and
  residual pieces.}
\end{table}

Two qualitative features stand out.
\begin{itemize}
  \item The zero-entry $(0,\pi,D)$-type seeds form a lower common-piece band,
    $z_0 \sim 18$--$19$ in MeV$^{1/2}$ units, already near Koide.
    The retained $(-\pi,D,X)$ continuations sit in a distinct higher band
    $z_0 \sim 29$--$40$.  Hence the pion sector does \emph{not} exhibit a
    single universal background charge: the common piece depends strongly on
    whether one is looking at a seed tuple or at its heavy continuation.
  \item The direct pion tuples share a common pattern: the pion supplies the
    dominant negative residual, the middle meson supplies a modest positive
    one, and the heaviest state supplies the large positive counterweight.
    In this sense the pion behaves like the exceptional light constituent
    around which the rest of the tuple organizes.
  \item The direct pion families and the retained inverse tuple do \emph{not}
    point to a single universal background scale: the direct families have
    $z_0 \sim 30$--$40$ in MeV$^{1/2}$ units, whereas the retained inverse
    tuple has $w_0 \sim 0.036$ in MeV$^{-1/2}$ units.
\end{itemize}

Taken together, the scan points to three robust regularities rather than a
large zoo of numerology: the factor-$3$ exact scaling along the quark chain,
the split pion common-piece bands separating the $(0,\pi,D)$ seeds from the
$(-\pi,D,X)$ continuations, and the coherence of the inverse down-type
family near $R_w=1$.  By contrast, the direct/inverse near-reciprocal hits
occur only sporadically, often with an additional overall minus sign, so the
present data do not support a single universal electric-magnetic or seesaw
duality between all direct and inverse preon charges.

Among the retained meson benchmarks of \S11, no kaon appears.
Moreover, the current benchmark set does not distinguish the neutral-kaon
basis $(K^0,\bar K^0)$ from the mass basis $(K_S,K_L)$.
If kaons are to participate in the same preon-pattern game, a dedicated
scan in those alternative kaon bases is still needed.

\end{document}
